\documentclass[11pt,a4paper]{article}

\usepackage[utf8]{inputenc}
\usepackage[T1]{fontenc}
\usepackage{lmodern}
\usepackage{microtype}
\usepackage[margin=2.6cm]{geometry}
\usepackage{amsmath,amssymb,amsfonts}
\usepackage{bm}
\usepackage{booktabs}
\usepackage{algorithm}
\usepackage{algpseudocode}
\usepackage{enumitem}
\usepackage[hidelinks]{hyperref}

\newcommand{\R}{\mathbb{R}}
\newcommand{\D}{\mathcal{D}}
\newcommand{\inner}[2]{\langle #1, #2 \rangle}
\newcommand{\norm}[1]{\lVert #1 \rVert}
\DeclareMathOperator*{\argmax}{arg\,max}
\DeclareMathOperator*{\argmin}{arg\,min}

\title{Sound decomposition strategies for creative hybridizations}
\author{Carmine-Emanuele Cella}
\date{}

\begin{document}
\maketitle

\begin{abstract}
A sound hybrid is a sound that belongs to two sounds at once: it takes its
form from one and its matter from the other. This paper presents hybridization
as a problem of representation: a sound is decomposed on a dictionary,
and the hybrid is obtained by reading the decomposition of one sound with the
elements of another. The kind of hybrid depends on the level at which the
representation operates, from the waveform to the spectrum to structured
decompositions. Two methods are described in detail. The first, Maple, is a
matching pursuit over parametric dictionaries (Gabor and gammatone atoms) and
over dictionaries made of recorded sounds, which approximates a target with
pieces of a corpus. The second, sound types, turns a sound into a hierarchy
of symbolic sequences (frames or notes, gestures, sections), generates new
sequences with $n$-grams, variable-order Markov models or a factor oracle, and
replaces the units of one sound with those of another at any level of the
hierarchy. The two meet when the units of sound types are made of the atoms
found by matching pursuit.
\end{abstract}

% =====================================================================
\section{Introduction}

\subsection{Hybridization}

Composers have always wished to make a sound that is two sounds at once: a
voice that is also a bell, an orchestra that speaks, a city that sings a
song. Such a sound is not a mixture. A mixture of two sounds is heard as two
sounds, because the auditory system groups the components of each source by
their common onsets, harmonicity and common fate~\cite{bregman1990}; the two
remain separable streams. A \emph{hybrid}, on the contrary, is heard as one
sound whose properties come from two origins. Wishart describes the
transformations that lead from one sound to another and the intermediate
states in which a sound is ambiguously both~\cite{wishart1994}; Smalley's
spectromorphology gives a vocabulary to the way the spectral content of a
sound (its matter) and its evolution in time (its form) can be thought of
separately~\cite{smalley1997}.

This separation is the key to hybridization. In most hybrids one sound gives
the \emph{form}: its temporal organisation, its rhythm, its envelope, its
articulation, sometimes its contour of brightness or pitch; the other gives
the \emph{matter}: the timbre, the spectral fine structure, the identity of
its sources. We call the first the \emph{target} and the second the
\emph{source} (or \emph{corpus}, when it is a collection of sounds). The
hybrid is the target rebuilt with the source's matter: John Cage's speech
made of a song by the Beach Boys, a voice played by a clarinet, a Bach
prelude played with the sound of a kitchen. The perceptual success of a
hybrid depends on a balance: if the target dominates, the source is not
recognised; if the source dominates, the target's form is lost. A good
hybrid lets the listener hear both: the speech and the song, the voice and
the clarinet.

\subsection{Decomposition and approximation}

To separate form and matter one needs a representation in which the two are
distinct objects. The general framework is \emph{decomposition on a
dictionary}. Let $x \in \R^N$ be a sound and let
$\D = \{ d_\gamma \}_{\gamma \in \Gamma}$ be a dictionary: a collection of
waveforms (\emph{atoms}), normalised so that $\norm{d_\gamma} = 1$, indexed by
a set of parameters $\Gamma$. A decomposition of $x$ on $\D$ is an
approximation
\begin{equation}
  x \;=\; \sum_{k=1}^{K} \alpha_k \, d_{\gamma_k} \;+\; r_K ,
  \label{eq:decomposition}
\end{equation}
where the $\alpha_k$ are coefficients, the $\gamma_k$ the atoms chosen, and
$r_K$ the residual: what the dictionary cannot (or does not need to)
explain. When $\D$ is a basis (the Fourier basis, a wavelet basis) the
decomposition is unique and exact; when $\D$ is \emph{redundant}, there are
infinitely many decompositions and one looks for a \emph{sparse} one: few
atoms that explain most of the energy. Redundancy is what makes a dictionary
expressive: Gabor's acoustic quanta~\cite{gabor1947}, grains of sound
localised in time and frequency, form a redundant dictionary in which a sound
is a cloud of events, as in granular synthesis~\cite{roads2001}.

The decomposition~\eqref{eq:decomposition} is a description of $x$ in the
terms of $\D$: the list $\{(\gamma_k, \alpha_k)\}$ (the \emph{book}) says
which elements of the dictionary sound when, and how strongly. Two remarks
make it a tool for hybridization.

\begin{enumerate}[leftmargin=*]
\item The book carries the \emph{form} of $x$ (positions, amplitudes, the
  order of events), while the atoms carry the \emph{matter}. If the
  dictionary is made of pieces of another sound, the approximation
  $\hat x = \sum \alpha_k d_{\gamma_k}$ is $x$ rebuilt with that sound's
  matter: a hybrid, whose degree of resemblance to $x$ is controlled by the
  number of atoms $K$ and by how well $\D$ can approximate $x$.
\item The book can be \emph{read with another dictionary}. Given a map
  $\phi : \D \to \D'$ that associates to each atom an atom of a second
  dictionary, the sound $\sum \alpha_k \, \phi(d_{\gamma_k})$ keeps the form
  of $x$ and takes the matter of $\D'$. The map can be a similarity in some
  descriptor space, or the identity of a symbolic label.
\end{enumerate}

In both cases the hybrid is an \emph{approximation of the target with a
specific dictionary}. Its quality is not the precision of the approximation
(an exact reconstruction is just the target) but the right amount of
\emph{infidelity}: the target must be recognisable, the dictionary audible.

\subsection{The level of the representation}

The kind of hybrid one obtains depends on the level at which the sounds are
represented, that is, on what the atoms are and how much structure they
carry.

\paragraph{Waveform.} At the level of samples the only operation is the
weighted sum $y = (1-\lambda)\, a + \lambda\, b$, possibly with a time-varying
$\lambda(t)$: a crossfade. It is the degenerate case of
\eqref{eq:decomposition} with the canonical basis of Dirac impulses, in which
form and matter cannot be separated. The result is a mixture or a
succession, never a hybrid.

\paragraph{Spectrum.} With the short-time Fourier transform
$X(m, k) = |X(m,k)|\, e^{i \varphi_X(m,k)}$ form and matter begin to separate.
\emph{Cross-synthesis} combines the magnitudes of one sound with the phases
of another, $Y = |A|\, e^{i\varphi_B}$, or more musically the spectral
envelope of one with the fine structure of the other,
$|Y| = E_A \cdot |B| / E_B$, where $E$ is a smoothed envelope (cepstral or
by linear prediction): this is the principle of the vocoder~\cite{dudley1939}
and of its digital descendants, the phase vocoder~\cite{flanagan1966,dolson1986}
and source-filter cross-synthesis by linear prediction~\cite{moorer1979}.
Here the target gives a slowly varying filter, the source an excitation. The
result is convincing when the target is a voice (its formants are its form)
and the source is rich and stationary; it is limited by the frame: the hybrid
is a pointwise product of two spectra at the same instant, and the temporal
organisation of the source is lost or imposed.

\paragraph{Decompositions with structured dictionaries.} When the atoms are
themselves sounds (grains, notes, strokes, gestures) the map
$\phi$ can operate between units that have duration, internal evolution and
identity. A voice approximated with clarinet notes is not a voice filtered by
a clarinet: it is a sequence of clarinet notes whose positions, amplitudes and
pitches follow the voice. A speech whose units are replaced, one by one, by
the units of a song that most resemble them is the speech made of the song,
with the song's textures, chords and voices recognisable inside the
articulation of the speech. The higher the level of the units, the more of
the source's own organisation survives: frames give a texture, notes a
re-orchestration, gestures and phrases a recomposition. The two methods of
this paper work at these levels: Maple at the level of atoms (parametric
or sampled), sound types at the level of symbols organised in a hierarchy.

% =====================================================================
\section{State of the art}

\paragraph{Spectral cross-synthesis and morphing.} Beyond the vocoder and
linear prediction~\cite{dudley1939,moorer1979,dolson1986}, sinusoidal and
spectral modelling~\cite{serra1990} made it possible to interpolate the
parameters of two sounds (frequencies, amplitudes, residual envelopes) and
to morph between them; automatic morphing aligns the two sounds in time and
in their spectral features before interpolating~\cite{slaney1996,caetano2011}.
Morphing produces intermediate sounds along a path between two endpoints; a
hybrid, as defined above, is rather a sound in which the two endpoints
coexist with different roles.

\paragraph{Concatenative synthesis and mosaicing.} Corpus-based concatenative
synthesis~\cite{schwarz2007} assembles a sound from units of a corpus chosen
by their descriptors. Musical mosaicing~\cite{zils2001} formulated the choice
of samples that follow a target sequence as a constraint problem; MATConcat
applied frame-based matching of descriptors to target-driven
synthesis~\cite{sturm2004}. CataRT~\cite{schwarz2006} made the approach
interactive: the units of a corpus are displayed in a descriptor space and
played by navigating it, or selected to follow a target in real time.
AudioGuide~\cite{hackbarth2013} uses a target to drive the selection and the
superposition of corpus segments, matching not only static descriptors but
their evolution in time, and lets several corpus segments overlap to
approximate a single target segment: a concatenative approach that is already
a decomposition, since the target is explained by a sum of sounds.
NMF-inspired mosaicing~\cite{driedger2015} keeps the activations of the
target while constraining the corpus to its own temporal continuity.

\paragraph{Dictionary-based methods.} Matching pursuit~\cite{mallat1993}
computes a sparse decomposition~\eqref{eq:decomposition} greedily, choosing
at each step the atom most correlated with the residual. Efficient
implementations for audio, with local updates and very large time-frequency
dictionaries, are provided by MPTK~\cite{krstulovic2006}. Sturm and
collaborators studied the analysis, visualisation and transformation of audio
through dictionary-based methods~\cite{sturm2009}, and cross-synthesis and
morphing by exchanging dictionaries or books between
sounds~\cite{collins2011}. Matching pursuit with dictionaries made of
recorded sounds turns concatenative synthesis into an approximation problem:
the corpus units are atoms, the target the signal to explain.

\paragraph{Target-based orchestration.} Computer-assisted orchestration can
be read as hybridization with an instrumental dictionary: a target sound is
approximated by a combination of instrumental notes, chosen among thousands
of samples of playing techniques, so that the orchestra imitates the
target. Orch\'id\'ee formulated it as multi-objective combinatorial
optimisation~\cite{carpentier2010}; Orchids extended it to dynamic targets.
Orchidea~\cite{cella2022orchidea,cella2018toolbox} solves target-based
dynamic orchestration with a modular framework, sparse and evolutionary
optimisation and large datasets of instrumental techniques such as
OrchideaSOL and TinySOL~\cite{cella2020orchideasol}. The orchestration of a
sound is a hybrid in which the form comes from the target and the matter
from the orchestra; the solutions are scores, i.e. symbolic books, that can
be performed.

\paragraph{Sound types.} The theory of sound types~\cite{cella2011soundtypes}
proposed to represent a sound by a sequence of symbols: the frames of the
sound are clustered in a feature space and each cluster, a \emph{sound
type}, is a symbol whose realisation is a set of sound fragments; the
temporal organisation of the sound becomes a sequence on a finite alphabet,
whose statistics ($n$-grams) can generate new sequences, hence new sounds
made of the original's matter. The approach was then applied to
hybridization~\cite{cella2013hybrid}: the types of a target are matched to
the closest types of a source in their respective feature spaces, so that the
target's sequence is played with the source's sounds, and the transition
probabilities of the two sounds can be combined to generate sequences that
belong to both.

\paragraph{Generation from symbolic sequences.} Variable-order Markov
models~\cite{ron1996,begleiter2004} adapt the length of the context to the
data; the factor oracle~\cite{allauzen1999} indexes all the factors of a
sequence in linear space and, navigated through its suffix links, generates
recombinations that preserve local continuity, as in machine
improvisation~\cite{assayag2004} and in the Variable Markov
Oracle~\cite{wang2014}. Surprise, the information content of an event given
its context, is a cue for the perception of
boundaries~\cite{abdallah2009,pearce2010}.

\paragraph{Neural methods.} Neural audio models learn latent spaces in which
timbres can be interpolated~\cite{engel2017}, or transfer the timbre of an
instrument to a performance by another through differentiable
synthesis~\cite{engel2020} or variational
autoencoders~\cite{caillon2021}. They offer powerful hybrids at the cost of
training on large corpora; the methods presented here require no training
and work with any pair of sounds, including very short and very unusual
ones.

% =====================================================================
\section{Methods}

Both methods operate on monophonic signals sampled at rate $f_s$. A
recurring tool is onset detection: the magnitude spectrogram $|X(m,k)|$
(frames of $1024$ samples, hop $256$) gives the spectral flux
\begin{equation}
  F(m) = \sum_k \max\bigl(0,\; |X(m,k)| - |X(m-1,k)|\bigr),
\end{equation}
whose local maxima above $\theta \cdot (\mathrm{med}_{21}F(m) + 0.01 \max F)$
are onsets~\cite{bello2005,dixon2006}, $\mathrm{med}_{21}$ being a running
median over 21 frames and $\theta$ a sensitivity; each onset is refined to
the sample where the energy rises in short blocks, and onsets closer than a
minimum gap are merged.

% ---------------------------------------------------------------------
\subsection{Maple: matching pursuit over parametric and sampled dictionaries}

Maple (\emph{matching pursuit linear expansion}) approximates a target with a
dictionary by matching pursuit~\cite{mallat1993}: starting from $r_0 = x$, at
each step it selects the atom with the largest projection on the residual,
\begin{equation}
  \gamma_k = \argmax_{\gamma \in \Gamma} \bigl|\inner{r_{k-1}}{d_\gamma}\bigr| ,
  \qquad
  \alpha_k = \inner{r_{k-1}}{d_{\gamma_k}} ,
  \qquad
  r_k = r_{k-1} - \alpha_k\, d_{\gamma_k},
\end{equation}
so that the energy of the residual decreases by $\alpha_k^2$ at each step,
$\norm{r_k}^2 = \norm{r_{k-1}}^2 - \alpha_k^2$. The book
$\{(\gamma_k, \alpha_k)\}$ is the description of the target; the
approximation after $K$ steps has a signal-to-residual ratio
$\mathrm{SRR}_K = 10 \log_{10} \bigl(\norm{x}^2 / \norm{r_K}^2\bigr)$.
Two families of dictionaries are used.

\subsubsection{Parametric dictionaries}

A parametric atom is a windowed sinusoid
\begin{equation}
  g_{s,u,\xi,\varphi}(t) = c \; w_s(t - u)\, \cos\bigl(2\pi \xi (t-u) + \varphi\bigr),
\end{equation}
with scale (window length) $s$, position $u$, frequency $\xi$ and phase
$\varphi$, and $c$ normalising it to unit energy. Two windows are offered.
The \emph{Gabor} window is a Gaussian truncated at $\pm 3\sigma$,
$w_s(t) = \exp\bigl(-18\,(t/s - \tfrac12)^2\bigr)$ for $0 \le t < s$, which
gives the best time-frequency concentration~\cite{gabor1947}. The
\emph{gammatone} window
$w_s(t) \propto (t/s)^3 \, e^{-12\, t/s}$ has a sharp attack and an
exponential decay: it fits percussive and resonant events, whose energy
precedes their tail, without the pre-echo of symmetric windows.

The scales are powers of two from $s_{\min}$ to $s_{\max}$ (by default 256
to 8192 samples, i.e.\ 6 to 186\,ms at 44.1\,kHz); the positions are on a
grid of a quarter of the scale, $u \in \{0, s/4, s/2, \dots\}$; the frequency
is continuous between $\max(f_{\min}, f_s/s)$ and $f_{\max}$; the phase is
continuous.

\paragraph{Selection with the STFT.} For each scale $s$ the residual is
analysed by an STFT with window $w_s$, hop $s/4$ and a transform of size
$M = 2s$ (zero-padded). For a frame at position $u$ and a bin $k$, with
$\omega_k = 2\pi k / M$, let
\begin{equation}
  a = \sum_t r(u+t)\, w_s(t) \cos(\omega_k t), \qquad
  b = \sum_t r(u+t)\, w_s(t) \sin(\omega_k t)
\end{equation}
be the correlations with the windowed cosine and sine. Because the windowed
cosine and sine are not orthogonal, the best atom of frequency $\omega_k$ at
$u$ is found by projecting on the plane they span. With the Gram terms
\begin{equation}
  C = \tfrac12 (W + c_2), \quad S = \tfrac12 (W - c_2), \quad X = \tfrac12 s_2, \qquad
  W = \sum_t w_s^2(t), \quad c_2 + i\, s_2 = \sum_t w_s^2(t)\, e^{i 2\omega_k t},
\end{equation}
the coefficients of the projection are
\begin{equation}
  \alpha = \frac{S a - X b}{C S - X^2}, \qquad
  \beta = \frac{C b - X a}{C S - X^2},
\end{equation}
and the energy captured is $E = \alpha a + \beta b$. The terms $c_2, s_2$ are
read from the transform of $w_s^2$ at bin $2k$, so that all bins of a frame
are evaluated with two FFTs. Each frame stores its best bin and energy.

\paragraph{Refinement.} The frame with the largest energy over all scales is
selected; the frequency is refined by a parabola through the logarithmic
magnitudes of the best bin and its neighbours, and the projection is
recomputed exactly at the refined frequency $\omega$. The atom is then
\begin{equation}
  \alpha \, w_s(t)\cos(\omega t) + \beta\, w_s(t)\sin(\omega t)
  = A \, w_s(t) \cos(\omega t - \varphi), \qquad
  A = \sqrt{\alpha^2 + \beta^2},\; \varphi = \operatorname{atan2}(\beta, \alpha),
\end{equation}
which is subtracted from the residual.

\paragraph{Local updates.} Only the frames that overlap the support of the
atom just subtracted change: for every scale, the frames intersecting
$[u, u+s)$ are re-analysed, and the others keep their stored maxima, as in
MPTK~\cite{krstulovic2006}. The cost of an iteration is therefore
proportional to the length of the atom, not to the length of the sound.

\paragraph{Stopping.} The pursuit stops after $K_{\max}$ atoms (2000) or when
the residual is $\mathrm{SRR}_{\text{target}}$ dB below the sound (30 dB).
Resynthesis is the sum of the atoms, optionally transposed by a ratio $\rho$
(each $\omega$ multiplied by $\rho$) or stretched in time by $\sigma$ (each
$u$ multiplied by $\sigma$, durations unchanged); because the atoms are
explicit, transposition and stretching are free of the artefacts of the
phase vocoder. Keeping only the fraction $q$ of the strongest atoms gives a
\emph{skeleton} of the sound.

\subsubsection{Sampled dictionaries}

A sampled dictionary is a collection of sound fragments taken from a corpus
of recordings. Each file can be cut in three ways:
\begin{itemize}[leftmargin=*]
  \item \emph{frames}: Hann-windowed segments of length $L$ (100\,ms) every
    hop $H$ (50 or 100\,ms);
  \item \emph{onsets}: the spans between consecutive onsets (at most
    $L_{\max}$ long), with 10\,ms fades;
  \item \emph{whole}: each file as one atom (at most $L_{\max}$ long).
\end{itemize}
Each fragment $d$ is normalised to unit energy and described by:
its prefix energies $P_d(m) = \sum_{t<m} d(t)^2$, $m = 0, \dots, n_d$; a
\emph{mel descriptor} $\mu_d \in \R^{40}$, the mean over its frames of the log
energies in 40 mel bands, with the mean removed and unit norm; its spectral
centroid.

\paragraph{Segments.} The target is cut into segments, either at its onsets
or in blocks of fixed length $B$ (a segment longer than $B$ is cut into
blocks too). Each segment $[a, a+S)$ is explained by $P$ \emph{components},
i.e.\ $P$ steps of pursuit restricted to atoms that start in the segment.

\paragraph{Preselection.} At each component, the residual segment is
described by its own mel descriptor $\mu_r$, and the $Q$ atoms of the
dictionary with the largest cosine similarity $\inner{\mu_r}{\mu_d}$ are kept
as candidates. This step makes large dictionaries (thousands of atoms)
tractable and acts as a perceptual prior: candidates are atoms whose spectral
envelope resembles the part of the target to explain.

\paragraph{Sliding by cross-correlation.} An atom of length $n$ can start at
any offset $\tau \in [0, S)$ of the segment, and it is matched over its whole
length, including the part that extends beyond the segment (up to the end of
the target). Let $W = \min(S + n, N - a)$ be the length of the residual
window considered. The correlation of the atom with the residual at every
offset,
\begin{equation}
  c(\tau) = \sum_{t=0}^{m(\tau) - 1} r(a + \tau + t)\, d(t), \qquad
  m(\tau) = \min(n, W - \tau),
\end{equation}
is computed for all $\tau$ at once with one FFT of size $\ge S + n$. Since
near the end of the signal only part of the atom fits, the gain of the
least-squares fit of the part that fits is
\begin{equation}
  G(\tau) = \frac{c(\tau)^2}{P_d\bigl(m(\tau)\bigr)}, \qquad
  \alpha(\tau) = \frac{c(\tau)}{P_d\bigl(m(\tau)\bigr)},
\end{equation}
discarding offsets where $P_d(m) < 0.05$ (too little of the atom fits). The
candidate and offset with the largest $G$ are chosen, the scaled atom is
subtracted from the residual over its whole length, and the book records its
index, start, coefficient, energy and the length of the part inside its
segment.

\paragraph{Synthesis.} The hybrid is
\begin{equation}
  \hat x(t) = \sum_k \alpha_k \, d_{\gamma_k}\!\bigl(\rho\,(t - \sigma u_k)\bigr),
\end{equation}
with transposition $\rho$ (by resampling) and stretch $\sigma$. With
\emph{tails} on, every atom plays whole and rings beyond its segment (as a
note that is allowed to resonate); with tails off it stops at the end of its
segment, which gives a drier, more articulated hybrid. A result that would
clip is scaled down.

\paragraph{Cross.} A book made with dictionary $\D_1$ can be read with a
dictionary $\D_2$: each atom $d \in \D_1$ is replaced by the atom of $\D_2$
whose mel descriptor is closest, with the same position and coefficient.
This is the map $\phi$ of the introduction: the target is explained once,
and its explanation is re-performed by another corpus.

\paragraph{Parameters and their effect.} Table~\ref{tab:maple} summarises
them. The number of components per segment is the main control of the
balance between target and source: one or two atoms per segment give a
hybrid in which the source is clearly heard and the target is a rhythm and a
contour; many atoms approach the target and the source dissolves into a
texture. Short frames and short blocks favour the target's spectral detail;
whole notes cut at onsets favour the source's identity (one recognisable
note per syllable). The number of candidates trades accuracy for speed.

\begin{table}[t]
\centering
\small
\begin{tabular}{lll}
\toprule
Parameter & Default & Role \\
\midrule
$s_{\min}, s_{\max}$ & 256, 8192 samples & scales of parametric atoms \\
$f_{\max}$ & 16 kHz & highest frequency of parametric atoms \\
$K_{\max}$, SRR & 2000, 30 dB & stopping (parametric) \\
cut, $L$, $H$, $L_{\max}$ & frames, 100, 100 ms, 1--2 s & how the corpus becomes atoms \\
segments, $B$, $\theta$ & onsets, 100 ms, 2 & how the target is cut \\
$P$ & 4 & atoms per segment \\
$Q$ & 16 & candidates after preselection \\
$\rho$, $\sigma$ & 1, 1 & transposition, time stretch \\
$q$ & 1 & share of the strongest atoms kept \\
tails & on & atoms ring whole or stop at their segment \\
\bottomrule
\end{tabular}
\caption{Parameters of Maple.}
\label{tab:maple}
\end{table}

\begin{algorithm}[t]
\caption{Matching pursuit with a sampled dictionary}
\label{alg:soundmp}
\begin{algorithmic}[1]
\Require target $x$, dictionary $\{d_j, P_{d_j}, \mu_{d_j}\}$, segments $\{[a_i, a_i+S_i)\}$, $P$, $Q$
\State $r \gets x$; book $\gets \emptyset$
\For{each segment $[a, a+S)$}
  \For{$p = 1, \dots, P$}
    \State $\mu_r \gets$ mel descriptor of $r[a, a+S)$
    \State $\mathcal{C} \gets$ the $Q$ atoms with largest $\inner{\mu_r}{\mu_{d_j}}$
    \For{$d \in \mathcal{C}$}
      \State $c(\tau) \gets$ correlation of $d$ with $r[a, a+S+n_d)$, $\tau \in [0, S)$ \Comment{FFT}
      \State $G(\tau) \gets c(\tau)^2 / P_d(m(\tau))$
    \EndFor
    \State $(d^*, \tau^*) \gets \argmax G$; $\alpha \gets c(\tau^*) / P_{d^*}(m(\tau^*))$
    \State $r[a+\tau^* + t] \gets r[a+\tau^*+t] - \alpha\, d^*(t)$ for $t < m(\tau^*)$
    \State append $(d^*, a+\tau^*, \alpha)$ to book
  \EndFor
\EndFor
\State \Return book, $r$
\end{algorithmic}
\end{algorithm}

% ---------------------------------------------------------------------
\subsection{Sound types: hierarchical symbolic representations}

The theory of sound types was introduced in~\cite{cella2011soundtypes} and
applied to hybridization in~\cite{cella2013hybrid}: a sound is represented by
a sequence of symbols, the \emph{types}, obtained by clustering its frames in
a feature space; the sequence can be modelled statistically and re-performed
with the sound's own fragments, or with the fragments of another sound. The
version presented here expands and improves that definition in five
directions: the units of the first level can be frames or notes (spans
between onsets); the representation is a hierarchy of levels (units,
gestures, sections), each level being a sequence of types; generation is done
at any level with $n$-grams, variable-order Markov models or a factor oracle,
and levels can be mixed; hybridization works frame by frame or unit by unit
at any level; and the units can be made of atoms of a matching pursuit.

\subsubsection{Features}

A sound is analysed by an STFT with Hann window of $N_w$ samples and hop
$h$ (by default 2048 and 512), the frames centred on multiples of $h$. Each
frame $m$ is described by a vector $\bm{f}_m$:
\begin{itemize}[leftmargin=*]
  \item 13 mel-frequency cepstral coefficients~\cite{davis1980}: the DCT of
    the log energies in 40 mel bands;
  \item the logarithm (base 2) of the spectral centroid, so that brightness
    is measured in octaves;
  \item the level in dB;
  \item optionally, the 12 pitch classes (chroma~\cite{fujishima1999}),
    computed from the spectral peaks (each peak, refined by parabolic
    interpolation, adds its magnitude to the nearest pitch class, partials 3,
    5, 6 and 7 of a lower peak being discounted), normalised to a maximum of
    1 and weighted by 3, so that harmony counts as much as timbre.
\end{itemize}
When the sound is given as a book of atoms (Section~\ref{sec:atoms}), the
frame features are the energies of the atoms present in each frame, in 40
logarithmic bands from 50\,Hz to 16\,kHz, compressed as $\log(1 + 10^3 e)$.

A stretch of frames $[m_0, m_1]$ of duration $T$ seconds is described by
\begin{equation}
  \bm{\delta}(m_0, m_1) = \bigl[\, \bm{\mu},\; \bm{\sigma},\; \log T \,\bigr],
  \label{eq:describe}
\end{equation}
the mean and standard deviation of its frame features and its log duration:
a description of the unit as a whole, comparable between units of different
lengths and between different sounds.

\subsubsection{Level 1: units}

Level 1 turns the sound into a sequence of \emph{units}, each with a type.
Two segmentations are possible.

\paragraph{Frames.} The frame features are standardised (each dimension to
zero mean and unit variance) and clustered into $k_1$ types by $k$-means
with $k$-means++ seeding~\cite{arthur2007}, keeping the best of four
starts. The labels are smoothed by a majority vote over a window of
$n_s$ frames (5), which removes isolated frames. A run of consecutive frames
of the same type is a unit. This segmentation is adapted to textures,
speech and sounds without clear onsets: the units are as long as the sound
stays in one region of the feature space.

\paragraph{Onsets.} The units are the spans between consecutive onsets
(a first span shorter than 50\,ms joins the next). Each unit is described as
a whole by~\eqref{eq:describe}; the descriptions are standardised and
clustered into $k_1$ types. This segmentation is adapted to instrumental
music: a unit is a note or a chord, and its type, with pitch classes in the
features, reflects both its timbre and its harmony.

In both cases each type $t$ has a centroid $\bar{\bm{\delta}}_t$, the mean
description of its units.

\subsubsection{Higher levels: gestures and sections}

The sequence of types $s_1, s_2, \dots, s_n$ of level $\ell$ is grouped into
the units of level $\ell+1$ (gestures, then sections) by cutting it where the
next type is \emph{surprising}. A bigram model with additive smoothing
estimates
\begin{equation}
  P(s_i = b \mid s_{i-1} = a) = \frac{\mathrm{count}(a, b) + \tfrac12}{\sum_{b'} \bigl(\mathrm{count}(a, b') + \tfrac12\bigr)},
\end{equation}
and the surprise of each unit is $h_i = -\log P(s_i \mid s_{i-1})$. With
$\bar h$ and $\sigma_h$ the mean and standard deviation of the surprises, a
boundary is placed before unit $i$ when
\begin{equation}
  h_i > \bar h + z\, \sigma_h \quad \text{and} \quad h_i \ge h_{i-1},\; h_i \ge h_{i+1},
\end{equation}
i.e.\ at a local peak of surprise above the threshold $z$ (1 by default), with
groups of at least 2 and at most 24 units. The choice follows the idea that
boundaries are perceived where expectation fails~\cite{abdallah2009,pearce2010}:
inside a gesture the next sound is predictable, at its end it is not.

Each group $G$ (units $i_0, \dots, i_1$ of level $\ell$) is described by
\eqref{eq:describe} over all its frames, augmented with the histogram of the
types of level $\ell$ it contains, weighted by duration,
\begin{equation}
  \bm{\eta}_G(t) = \frac{\sum_{i \in G,\, s_i = t} T_i}{\sum_{i \in G} T_i},
\end{equation}
so that two gestures are similar when they sound alike on average
\emph{and} are made of similar sounds. The augmented descriptions are
standardised and clustered into $k_{\ell+1}$ types. The procedure is applied
recursively; three levels are used in practice ($k = 12, 6, 4$ by default;
for a piano prelude cut at its onsets, $16, 8, 4$ yield notes, phrases and
sections). Each unit at every level keeps its span in samples and the range
of units of the level below it covers.

\subsubsection{Generation}

Generation produces a new sequence of units of a chosen level, until a
requested duration, and realises it with the original's sound. Three models
are available.

\paragraph{$n$-grams.} The classic model of~\cite{cella2011soundtypes}: the
next type is drawn from the distribution of the types that followed the last
$n-1$ types in the original. When the context was never seen, the
generation restarts from a random place, copying $n-1$ original units.
Small $n$ wanders; large $n$ copies long stretches, since few contexts are
shared.

\paragraph{Variable-order Markov.} The next type is drawn from the
distribution that follows the longest context, of length at most $D$, that
appears at least $c_{\min}$ times in the original, backing off to shorter
contexts (down to the unconditional distribution) when longer ones are
rare~\cite{ron1996,begleiter2004}. This adapts the memory to the material:
long where the music repeats, short where it does not.

\paragraph{Factor oracle.} The factor oracle~\cite{allauzen1999} of the
sequence $s_1 \dots s_n$ is an automaton with states $0, \dots, n$, a
transition $i-1 \to i$ labelled $s_i$, additional forward transitions that
recognise all factors of the sequence, and a suffix link $\mathrm{sfx}(i)$
from each state to the state where the longest repeated suffix of
$s_1 \dots s_i$ ends. The length of that repeated suffix,
$\mathrm{lrs}(i)$, is computed during construction~\cite{lefebvre2000}.
Navigation~\cite{assayag2004,wang2014} starts at state 0; at state $i$:
\begin{itemize}[leftmargin=*]
  \item with probability $\kappa$ (the \emph{continuity}), it moves to $i+1$,
    i.e.\ plays the next unit of the original;
  \item otherwise it jumps: the candidates are $\mathrm{sfx}(i)$ and every
    state $j \ne i$ with $\mathrm{sfx}(j) = \mathrm{sfx}(i)$, provided their
    repeated suffix is at least $\lambda$ units long
    ($\mathrm{lrs} \ge \lambda$, the \emph{context}); one is chosen at random
    and the next unit after it is played.
\end{itemize}
At the end of the sequence the navigation restarts. With $\kappa = 1$ the
oracle replays the original; as $\kappa$ decreases it recombines it, always
jumping between places where the last $\lambda$ types are the same, so that
every junction is locally plausible. Unlike $n$-grams, the oracle generates
unit \emph{indices}, not types: it keeps track of where it is in the
original, and long stretches of the original survive between jumps.

\paragraph{Realisation.} For the $n$-gram and Markov models a type is
realised by one of its units: preferably the unit that follows the last one
played in the original (so that original continuations are kept whenever the
model allows them), otherwise a random unit of that type. The chosen units
are concatenated with crossfades of $x_f$ ms (10 by default). To avoid holes
and clicks, each unit is taken with $x_f/2$ ms of its original context on
both sides, which become the material of the fades; with $\kappa = 1$ the
concatenation reproduces the original almost exactly.

\paragraph{Mixing levels.} A sequence generated at level $\ell \ge 2$
decides the \emph{form} (which gestures or sections, in which order); the
content of each of its units can itself be regenerated at the level below.
For a unit $U$ of level $\ell$ covering units $f, \dots, l$ of level $\ell-1$,
the realisation keeps its first sub-unit $f$ and its number of sub-units
$l - f + 1$, and replaces the rest by an oracle walk on level $\ell-1$ that
starts at $f$ with an inner continuity $\kappa'$:
\begin{equation}
  f \;\to\; \mathrm{step}_{\kappa'}(f) \;\to\; \mathrm{step}_{\kappa'}^2(f) \;\to\; \dots
  \qquad (l - f + 1 \text{ sub-units}),
\end{equation}
recursively down to level 1. With $\kappa' = 1$ the units are played as
they are; with $\kappa' < 1$ the phrases still begin as the original ones
and last about as long, but their notes are new. Mixing levels gives a
family of results between two extremes: recombining notes (level 1), which
creates a new surface without phrasing of its own, and recombining phrases
or sections (levels 2 and 3), which keeps the inside of the units intact and
changes the form. The mixture keeps the large form of one level and the
detail of another: for instance the plan of sections of a prelude, with
every phrase rewritten.

\begin{algorithm}[t]
\caption{Factor oracle step and level mixing}
\label{alg:oracle}
\begin{algorithmic}[1]
\Function{Step}{$i$, $\kappa$, $\lambda$}
  \If{$i < n$ \textbf{and} ($\mathrm{rand}() < \kappa$ \textbf{or} $i = 0$)} \Return $i+1$ \EndIf
  \State $\mathcal{J} \gets \{\mathrm{sfx}(i)\} \cup \{ j \neq i : \mathrm{sfx}(j) = \mathrm{sfx}(i) \}$, keeping those with $\mathrm{lrs} \ge \lambda$
  \State $j \gets$ random element of $\mathcal{J}$ (a random state if empty)
  \State \Return $\min(j+1, n)$
\EndFunction
\Function{Nest}{sequence $u_1 \dots u_m$ at level $\ell$, $\kappa'$}
  \For{$L = \ell, \ell-1, \dots, 2$}
    \State $v \gets ()$
    \For{each $u$ in the sequence}
      \State $(f, l) \gets$ sub-units of $u$ at level $L-1$; append $f$ to $v$; $i \gets f+1$
      \For{$q = 1, \dots, l - f$} $i \gets \textsc{Step}(i, \kappa', \lambda)$ on level $L-1$; append $i-1$ to $v$ \EndFor
    \EndFor
    \State sequence $\gets v$
  \EndFor
  \State \Return sequence (units of level 1)
\EndFunction
\end{algorithmic}
\end{algorithm}

\subsubsection{Hybridization: rematch}

Hybridization replaces the units of a target $A$ with units of a source
$B$, both analysed into types. The central choice is the space in which they
are compared. Each sound's features are standardised \emph{in its own
space}: $\bm{z}^A = (\bm{f}^A - \bm{m}_A)/\bm{s}_A$ and
$\bm{z}^B = (\bm{f}^B - \bm{m}_B)/\bm{s}_B$, with the mean and standard
deviation of each sound. Distances are then relative: the brightest part of
$A$ corresponds to the brightest part of $B$, the loudest to the loudest, so
that the whole range of the source is used to render the variations of the
target, whatever the absolute difference between the two (a dry speech and
a full mix). A joint standardisation, measuring absolute distances, is also
possible; it tends to map most of the target onto the few source sounds that
are absolutely closest to it.

\paragraph{Frame by frame (level 0).} This is the rematch of the original
theory~\cite{cella2013hybrid}. Let $\bar{\bm{z}}^B_t$ be the centroid of the
frames of $B$ of type $t$ (level 1). For each frame $i$ of $A$:
\begin{equation}
  t^*(i) = \argmin_t \norm{\bm{z}^A_i - \bar{\bm{z}}^B_t}, \qquad
  j^*(i) = \argmin_{j : \mathrm{type}(j) = t^*} \norm{\bm{z}^A_i - \bm{z}^B_j},
\end{equation}
the closest type of $B$, then its closest frame. The frame $j^*$ of $B$ is
taken as a grain: the $N_B$ samples centred on it, with a Hann window $v$.
The grains are overlap-added at the frames of $A$ and normalised by the sum
of the windows,
\begin{equation}
  y(t) = \frac{\sum_i g_i \, v(t - i h_A + N_B/2)\; b\bigl(t - i h_A + j^*(i)\, h_B\bigr)}{\max\Bigl(\sum_i v(t - i h_A + N_B/2),\; \tfrac14 V\Bigr)},
  \qquad
  g_i = \min\Bigl(10,\; \bigl(e^A_i / e^B_{j^*(i)}\bigr)^{c}\Bigr),
\end{equation}
where $V$ is the maximum of the window sum (the floor avoids boosting the
edges), $e^A_i$ and $e^B_j$ are the frames' amplitudes and $c \in [0,1]$ is
the \emph{contour}: with $c = 0$ the grains keep their own level, with
$c = 1$ they take the target's. With long grains ($N_B = 8192$ samples,
$h = 2048$, i.e.\ 186\,ms grains overlapping four times) the result is a
continuous texture of the source that follows the target's timbre and
loudness frame by frame: the speech is heard through the song.

\paragraph{Unit by unit (levels $\ge 1$).} At level $\ell$, each unit
$u$ of $A$ is described by~\eqref{eq:describe}; the descriptions of the units
of $A$ and of $B$ are standardised in their own spaces; the type of $B$ with
the closest centroid is chosen, then its unit closest to $u$ (the log
duration being part of the description, the closest unit is also close in
length). The units of $B$ are placed either at the times of $A$, cut or
continued to $A$'s durations (\emph{fit}), or one after the other at their
own lengths, and joined by crossfades. At level 1 a speech becomes a
sequence of song fragments, one per syllable or sound; at level 2 each
gesture of the speech becomes a whole gesture of the song, and the source's
own phrases appear within the target's phrasing.

\paragraph{Contour.} At levels $\ge 1$ the target's loudness can be
imposed afterwards: with $\mathrm{env}$ the RMS over 20\,ms windows,
\begin{equation}
  y'(t) = y(t) \cdot \min\Bigl(10,\; \Bigl(\frac{\mathrm{env}_A(t)}{\max(\mathrm{env}_y(t), 10^{-3} \max \mathrm{env}_y)}\Bigr)^{c}\Bigr),
\end{equation}
then $y'$ is rescaled to the RMS of $y$: the shape of the target's dynamics
at the level of the source. The limit of $+20$\,dB and the floor on
$\mathrm{env}_y$ prevent the noise floor of the source from being pulled up
in the target's silences. The contour is what makes the target's rhythm and
phrasing intelligible in the hybrid.

\begin{table}[t]
\centering
\small
\begin{tabular}{lll}
\toprule
Parameter & Default & Role \\
\midrule
$N_w$, $h$ & 2048, 512 & analysis window and hop (8192, 2048 for level-0 grains) \\
segmentation & frames & level 1 from runs of frame types, or from onsets \\
pitch & off & add the 12 pitch classes to the features \\
$k_1, k_2, k_3$ & 12, 6, 4 & types per level \\
levels & 2 & depth of the hierarchy (1 to 3) \\
$n_s$ & 5 & majority smoothing of frame labels \\
$z$ & 1 & surprise threshold for boundaries \\
model & oracle & $n$-grams ($n$), Markov ($D$, $c_{\min}$), oracle ($\kappa$, $\lambda$) \\
$\kappa'$ & 1 & inner continuity when mixing levels \\
$x_f$ & 10 ms & crossfades between units \\
level (rematch) & 1 & 0 frames, 1 units, 2 gestures \\
fit, contour $c$, space & on, 0, own & timing, loudness, standardisation \\
\bottomrule
\end{tabular}
\caption{Parameters of sound types.}
\label{tab:st}
\end{table}

\subsubsection{Sound types made of atoms}
\label{sec:atoms}

The two methods connect when the sound analysed into types is not a
recording but a book of atoms. Given a book $\{(d_k, u_k, \alpha_k)\}$ of
the target $x$ on a dictionary $\D$, the frame features are the energies of
the atoms active in each frame, in 40 logarithmic bands; an atom whose
support covers several frames contributes to each of them. Types, units,
gestures and sections are then computed as before, and a unit spanning
$[a, b)$ is realised by resynthesising the atoms whose start falls in it:
\begin{equation}
  \hat x_{[a,b)}(t) = \sum_{k :\, a \le u_k < b} \alpha_k\, d_k(t - u_k + a).
\end{equation}
This has three consequences.
\begin{enumerate}[leftmargin=*]
  \item A book made with a sampled dictionary is already a hybrid; its types
    describe the hybrid, and generation recombines it: a voice explained by
    clarinet notes becomes a new clarinet piece with the voice's phrasing,
    reordered by the oracle.
  \item Units rebuilt from atoms have no cuts: the atoms of a unit ring on
    into the next one, and the junctions are as smooth as the atoms
    themselves. Transposition, time stretching and the selection of the
    strongest atoms apply to the generated sequences too.
  \item Two books made with different dictionaries but described in the same
    band space can be rematched: the symbolic level and the atomic level
    become two stages of the same hybridization, the first deciding
    \emph{what} replaces what, the second \emph{with which matter}.
\end{enumerate}
The chain is therefore: decompose the target with the atoms of a source
(Maple), describe the decomposition as types at several levels (sound
types), generate or rematch at the level that suits the musical intention,
and resynthesise from atoms.

% =====================================================================
\section{Conclusion}

The two methods presented here treat hybridization as the approximation of
one sound with the elements of another, at levels of representation that
carry more structure than the spectrum: atoms (parametric or sampled) and
symbolic types organised in a hierarchy. Their creative applications are
many: rebuilding a voice with an instrument (or an instrument with the
sounds of a place), at a controllable degree of resemblance; making the
skeleton of a sound audible by keeping its strongest atoms; transposing and
stretching sounds without spectral artefacts; recomposing a piece from its
own notes, phrases or sections, with models that range from free
recombination to near-literal continuation; mixing the form of one level
with the detail of another; playing the articulation of a speech with the
textures and voices of a song, frame by frame or gesture by gesture. Both
methods are fast enough to be used interactively on sounds of a few minutes
and require no training, so they work on any material, including the
composer's own recordings.

This paper describes the algorithms and their parameters; it contains no
testing or evaluation. Further development would need objective measures of
the trade-off between the target and the source in a hybrid (for instance
the correlation of their envelopes and of their spectral features with the
result, and the recognisability of each), listening tests on the perceived
fusion and identity, a study of the effect of the parameters, and
comparisons with concatenative and neural approaches. Natural extensions are
dictionaries learned from the corpus rather than cut from it, joint models
of the type sequences of two sounds for generation that belongs to both, and
sequence models of larger capacity (such as transformers) trained on the
symbolic sequences produced by sound types.

% =====================================================================
\begin{thebibliography}{99}
\small

\bibitem{abdallah2009}
S.~Abdallah and M.~Plumbley.
\newblock Information dynamics: patterns of expectation and surprise in the
  perception of music.
\newblock \emph{Connection Science}, 21(2--3):89--117, 2009.

\bibitem{allauzen1999}
C.~Allauzen, M.~Crochemore, and M.~Raffinot.
\newblock Factor oracle: a new structure for pattern matching.
\newblock In \emph{SOFSEM'99: Theory and Practice of Informatics}, LNCS 1725,
  pages 295--310. Springer, 1999.

\bibitem{arthur2007}
D.~Arthur and S.~Vassilvitskii.
\newblock k-means++: the advantages of careful seeding.
\newblock In \emph{Proceedings of the ACM-SIAM Symposium on Discrete
  Algorithms (SODA)}, pages 1027--1035, 2007.

\bibitem{assayag2004}
G.~Assayag and S.~Dubnov.
\newblock Using factor oracles for machine improvisation.
\newblock \emph{Soft Computing}, 8(9):604--610, 2004.

\bibitem{begleiter2004}
R.~Begleiter, R.~El-Yaniv, and G.~Yona.
\newblock On prediction using variable order Markov models.
\newblock \emph{Journal of Artificial Intelligence Research}, 22:385--421,
  2004.

\bibitem{bello2005}
J.~P. Bello, L.~Daudet, S.~Abdallah, C.~Duxbury, M.~Davies, and M.~Sandler.
\newblock A tutorial on onset detection in music signals.
\newblock \emph{IEEE Transactions on Speech and Audio Processing},
  13(5):1035--1047, 2005.

\bibitem{bregman1990}
A.~S. Bregman.
\newblock \emph{Auditory Scene Analysis: The Perceptual Organization of
  Sound}.
\newblock MIT Press, 1990.

\bibitem{caetano2011}
M.~Caetano and X.~Rodet.
\newblock Sound morphing by feature interpolation.
\newblock In \emph{Proceedings of ICASSP}, 2011.

\bibitem{caillon2021}
A.~Caillon and P.~Esling.
\newblock RAVE: a variational autoencoder for fast and high-quality neural
  audio synthesis.
\newblock arXiv:2111.05011, 2021.

\bibitem{carpentier2010}
G.~Carpentier, G.~Assayag, and E.~Saint-James.
\newblock Solving the musical orchestration problem using multiobjective
  constrained optimization with a genetic local search approach.
\newblock \emph{Journal of Heuristics}, 16(5):681--714, 2010.

\bibitem{cella2011soundtypes}
C.-E. Cella.
\newblock Sound-types: a new framework for symbolic sound analysis and
  synthesis.
\newblock In \emph{Proceedings of the International Computer Music
  Conference (ICMC)}, Huddersfield, 2011.

\bibitem{cella2013hybrid}
C.-E. Cella and J.~J. Burred.
\newblock Advanced sound hybridizations by means of the theory of
  sound-types.
\newblock In \emph{Proceedings of the International Computer Music
  Conference (ICMC)}, Perth, 2013.

\bibitem{cella2018toolbox}
C.-E. Cella and P.~Esling.
\newblock Open-source modular toolbox for computer-aided orchestration.
\newblock In \emph{Proceedings of the International Conference on
  Technologies for Music Notation and Representation (TENOR)}, 2018.

\bibitem{cella2020orchideasol}
C.-E. Cella, D.~Ghisi, V.~Lostanlen, F.~L\'evy, J.~Fineberg, and Y.~Maresz.
\newblock OrchideaSOL: a dataset of extended instrumental techniques for
  computer-aided orchestration.
\newblock In \emph{Proceedings of the International Computer Music
  Conference (ICMC)}, 2020.

\bibitem{cella2022orchidea}
C.-E. Cella.
\newblock Orchidea: a comprehensive framework for target-based
  computer-assisted dynamic orchestration.
\newblock \emph{Journal of New Music Research}, 2022.

\bibitem{collins2011}
N.~Collins and B.~L. Sturm.
\newblock Sound cross-synthesis and morphing using dictionary-based methods.
\newblock In \emph{Proceedings of the International Computer Music
  Conference (ICMC)}, 2011.

\bibitem{davis1980}
S.~Davis and P.~Mermelstein.
\newblock Comparison of parametric representations for monosyllabic word
  recognition in continuously spoken sentences.
\newblock \emph{IEEE Transactions on Acoustics, Speech, and Signal
  Processing}, 28(4):357--366, 1980.

\bibitem{dixon2006}
S.~Dixon.
\newblock Onset detection revisited.
\newblock In \emph{Proceedings of the International Conference on Digital
  Audio Effects (DAFx)}, 2006.

\bibitem{dolson1986}
M.~Dolson.
\newblock The phase vocoder: a tutorial.
\newblock \emph{Computer Music Journal}, 10(4):14--27, 1986.

\bibitem{driedger2015}
J.~Driedger, T.~Pr\"atzlich, and M.~M\"uller.
\newblock Let it Bee -- towards NMF-inspired audio mosaicing.
\newblock In \emph{Proceedings of the International Society for Music
  Information Retrieval Conference (ISMIR)}, 2015.

\bibitem{dudley1939}
H.~Dudley.
\newblock Remaking speech.
\newblock \emph{Journal of the Acoustical Society of America},
  11(2):169--177, 1939.

\bibitem{engel2017}
J.~Engel, C.~Resnick, A.~Roberts, S.~Dieleman, D.~Eck, K.~Simonyan, and
  M.~Norouzi.
\newblock Neural audio synthesis of musical notes with WaveNet
  autoencoders.
\newblock In \emph{Proceedings of the International Conference on Machine
  Learning (ICML)}, 2017.

\bibitem{engel2020}
J.~Engel, L.~Hantrakul, C.~Gu, and A.~Roberts.
\newblock DDSP: differentiable digital signal processing.
\newblock In \emph{Proceedings of the International Conference on Learning
  Representations (ICLR)}, 2020.

\bibitem{flanagan1966}
J.~L. Flanagan and R.~M. Golden.
\newblock Phase vocoder.
\newblock \emph{Bell System Technical Journal}, 45(9):1493--1509, 1966.

\bibitem{fujishima1999}
T.~Fujishima.
\newblock Realtime chord recognition of musical sound: a system using Common
  Lisp Music.
\newblock In \emph{Proceedings of the International Computer Music
  Conference (ICMC)}, 1999.

\bibitem{gabor1947}
D.~Gabor.
\newblock Acoustical quanta and the theory of hearing.
\newblock \emph{Nature}, 159:591--594, 1947.

\bibitem{hackbarth2013}
B.~Hackbarth, N.~Schnell, P.~Esling, and D.~Schwarz.
\newblock Composing morphology: concatenative synthesis as an intuitive
  medium for prescribing sound in time.
\newblock \emph{Contemporary Music Review}, 32(1):49--59, 2013.

\bibitem{krstulovic2006}
S.~Krstulovi\'c and R.~Gribonval.
\newblock MPTK: matching pursuit made tractable.
\newblock In \emph{Proceedings of ICASSP}, 2006.

\bibitem{lefebvre2000}
A.~Lefebvre and T.~Lecroq.
\newblock Computing repeated factors with a factor oracle.
\newblock In \emph{Proceedings of the Australasian Workshop on Combinatorial
  Algorithms (AWOCA)}, pages 145--158, 2000.

\bibitem{mallat1993}
S.~Mallat and Z.~Zhang.
\newblock Matching pursuits with time-frequency dictionaries.
\newblock \emph{IEEE Transactions on Signal Processing}, 41(12):3397--3415,
  1993.

\bibitem{moorer1979}
J.~A. Moorer.
\newblock The use of linear prediction of speech in computer music
  applications.
\newblock \emph{Journal of the Audio Engineering Society}, 27(3):134--140,
  1979.

\bibitem{pearce2010}
M.~T. Pearce, D.~M\"ullensiefen, and G.~A. Wiggins.
\newblock The role of expectation and probabilistic learning in auditory
  boundary perception: a model comparison.
\newblock \emph{Perception}, 39(10):1365--1389, 2010.

\bibitem{roads2001}
C.~Roads.
\newblock \emph{Microsound}.
\newblock MIT Press, 2001.

\bibitem{ron1996}
D.~Ron, Y.~Singer, and N.~Tishby.
\newblock The power of amnesia: learning probabilistic automata with variable
  memory length.
\newblock \emph{Machine Learning}, 25(2--3):117--149, 1996.

\bibitem{schwarz2006}
D.~Schwarz, G.~Beller, B.~Verbrugghe, and S.~Britton.
\newblock Real-time corpus-based concatenative synthesis with CataRT.
\newblock In \emph{Proceedings of the International Conference on Digital
  Audio Effects (DAFx)}, 2006.

\bibitem{schwarz2007}
D.~Schwarz.
\newblock Corpus-based concatenative synthesis.
\newblock \emph{IEEE Signal Processing Magazine}, 24(2):92--104, 2007.

\bibitem{serra1990}
X.~Serra and J.~O. Smith.
\newblock Spectral modeling synthesis: a sound analysis/synthesis system
  based on a deterministic plus stochastic decomposition.
\newblock \emph{Computer Music Journal}, 14(4):12--24, 1990.

\bibitem{slaney1996}
M.~Slaney, M.~Covell, and B.~Lassiter.
\newblock Automatic audio morphing.
\newblock In \emph{Proceedings of ICASSP}, 1996.

\bibitem{smalley1997}
D.~Smalley.
\newblock Spectromorphology: explaining sound-shapes.
\newblock \emph{Organised Sound}, 2(2):107--126, 1997.

\bibitem{sturm2004}
B.~L. Sturm.
\newblock MATConcat: an application for exploring concatenative sound
  synthesis using MATLAB.
\newblock In \emph{Proceedings of the International Computer Music
  Conference (ICMC)}, 2004.

\bibitem{sturm2009}
B.~L. Sturm, C.~Roads, A.~McLeran, and J.~J. Shynk.
\newblock Analysis, visualization, and transformation of audio signals using
  dictionary-based methods.
\newblock \emph{Journal of New Music Research}, 38(4):325--341, 2009.

\bibitem{wang2014}
C.~Wang and S.~Dubnov.
\newblock Guided music synthesis with variable Markov oracle.
\newblock In \emph{Proceedings of the Musical Metacreation Workshop (AIIDE)},
  2014.

\bibitem{wishart1994}
T.~Wishart.
\newblock \emph{Audible Design: A Plain and Easy Introduction to Practical
  Sound Composition}.
\newblock Orpheus the Pantomime, 1994.

\bibitem{zils2001}
A.~Zils and F.~Pachet.
\newblock Musical mosaicing.
\newblock In \emph{Proceedings of the International Conference on Digital
  Audio Effects (DAFx)}, 2001.

\end{thebibliography}

\end{document}
